\subsection{Prefetching}

With the transition to a 4-way instruction fetch, we observed a notable drop in instruction cache (i-cache) hit rates across several benchmark programs. To address this, we implemented a prefetching mechanism for the i-cache. Specifically, while serving pending fetch requests, we preemptively fetch subsequent instruction memory blocks into a stream buffer. This technique helps maintain a steady flow of useful requests to memory, thereby better utilizing available memory bandwidth.

The size of the stream buffer plays a critical role in overall CPI, as illustrated in Figure~\ref{PrefetchBufferSize}. Increasing the buffer size from 2 to 4 entries resulted in up to a 50\% reduction in CPI for certain workloads.

\begin{figure}[htbp]
\centerline{\includegraphics[width=0.8\columnwidth]{Analysis/PrefetchBufferSize}}
\caption{Impact of Prefetch Buffer Size on CPI. Blue bars correspond to a size of 2; red line indicates results for size 4.}
\label{PrefetchBufferSize}
\end{figure}

From the hit rate data in Figure~\ref{ICacheHitRate}, it is evident that the stream buffer size is a significant bottleneck in the fetch stage. Given constraints on buffer capacity, one potential optimization is to more aggressively promote data from the stream buffer into the i-cache, improving the likelihood of hits and reducing fetch stalls.

\begin{figure}[htbp]
\centerline{\includegraphics[width=0.7\columnwidth]{Analysis/ICacheHitRate}}
\caption{Instruction Cache Hit Rate vs. Stream Buffer Size}
\label{ICacheHitRate}
\end{figure}

\subsection{Banked Instruction Cache}

To support 4-way instruction fetch, we often need to access two instruction memory blocks concurrently. However, our original i-cache design featured only a single read port, making such parallel access infeasible. To resolve this, we modified the i-cache to include two independent banks, enabling simultaneous dual-block access. This architectural change led to an average CPI reduction of approximately 1\%.

\subsection{Issuing with Age-Based Logic}

In the baseline design, the issue logic selects instructions for execution in a pseudo-random fashion. Ideally, the processor should issue instructions whose execution will unlock the greatest number of dependent instructions. While tracking all dependencies precisely would be prohibitively complex, a useful heuristic is to prioritize older instructions, as they are more likely to have dependent operations waiting in the pipeline.

This strategy becomes particularly relevant when the number of ready reservation station entries exceeds the available functional unit or CDB capacity. To quantify this, we define the issue rate as:

$$
\frac{\text{\#entries\_issued}}{\text{\#entries\_ready}}
$$

Figure~\ref{IssueRate} shows the issue rate across various benchmarks.

\begin{figure}[htbp]
\centerline{\includegraphics[width=0.6\columnwidth]{Analysis/IssueRate}}
\caption{Issue Rate Across Benchmarks}
\label{IssueRate}
\end{figure}

In several programs, the issue rate hovers around 50\%, suggesting that age-based issuing has the potential to improve performance by making better use of available execution resources.

The performance impact of this heuristic is shown in Figure~\ref{AgeLogic}.

\begin{figure}[htbp]
\centerline{\includegraphics[width=0.7\columnwidth]{Design/AgeLogic}}
\caption{Performance Impact of Age-Based Issuing}
\label{AgeLogic}
\end{figure}

On average, we observed a modest 1.5\% improvement in performance. However, some programs experienced a noticeable performance degradation. This suggests that in certain scenarios, issuing younger instructions can unlock more parallelism than strictly following program order. Furthermore, the added complexity of age-based logic introduced non-trivial latency to the issue stage. Based on these trade-offs, we ultimately chose to revert to the simpler, baseline issue logic.
